\section{Station firmware}
\label{sec:firmware-station}

The station runs on ESP-IDF. It combines a NimBLE BLE central with a Wi-Fi client, so a single
ESP32-S3 both listens to the collar and talks to the cloud. It is built and flashed with
\texttt{idf.py} over USB serial.

\subsection{Responsibilities}

\begin{enumerate}[leftmargin=1.5em]
  \item Scan for registered collars and read their telemetry advertisements.
  \item In production mode, relay the collar's classification to the cloud.
  \item In training mode, connect and capture short audio clips (with time-aligned motion) and
        upload them for labelling and training.
  \item Gate, download, and push new models to the collar over the air.
\end{enumerate}

\subsection{Capture and relay}

When relaying, the station maps the collar's compact telemetry (mode, class index, confidence)
onto the cat's record in the database, writing both a live ``current'' state and discrete
activity events as episodes begin and end. Events carry the model version and class index so the
dashboard can resolve them to human-readable labels even as the label set changes
(\cref{sec:cloud-backend}).

When capturing, the station records audio back-to-back while a registered collar is within a
configurable signal-strength gate, uploads each clip to Cloud Storage, and records its metadata
in the database for the dashboard's labelling UI. A manual capture path ignores the range gate,
which is used to collect behaviours such as purring that happen while the cat rests rather than
at the bowl.

\subsection{Robustness: scanning after OTA}

The station uses one BLE radio for several jobs, and OTA push exposed a subtle defect. To push a
model the station must cancel scanning and open a connection to the collar; the OTA path used its
own GAP callback and never restarted scanning when it finished, so after an update the station
went deaf and the dashboard reported the collar offline. The fix is an idempotent
\texttt{ble\_resume\_scan()} (guarded on the discovery-active flag) that the OTA path calls on
both completion and bail-out. This was verified on hardware.

\subsection{Configuration}

Wi-Fi credentials are provisioned to the station over USB during device pairing in the
dashboard and are never stored in the repository (see the \texttt{.gitignore} note on
\texttt{secrets.h}). The capture range (signal threshold and dwell) and the operating mode are
configured per station from the dashboard and read by the firmware from the database.
